% ----------------------------------------------------
%
%          Denotational semantics
%
% ----------------------------------------------------

%% ODER: format ==         = "\mathrel{==}"
%% ODER: format /=         = "\neq "
%
%
\makeatletter
\@ifundefined{lhs2tex.lhs2tex.sty.read}%
  {\@namedef{lhs2tex.lhs2tex.sty.read}{}%
   \newcommand\SkipToFmtEnd{}%
   \newcommand\EndFmtInput{}%
   \long\def\SkipToFmtEnd#1\EndFmtInput{}%
  }\SkipToFmtEnd

\newcommand\ReadOnlyOnce[1]{\@ifundefined{#1}{\@namedef{#1}{}}\SkipToFmtEnd}
\usepackage{amstext}
\usepackage{amssymb}
\usepackage{stmaryrd}
\DeclareFontFamily{OT1}{cmtex}{}
\DeclareFontShape{OT1}{cmtex}{m}{n}
  {<5><6><7><8>cmtex8
   <9>cmtex9
   <10><10.95><12><14.4><17.28><20.74><24.88>cmtex10}{}
\DeclareFontShape{OT1}{cmtex}{m}{it}
  {<-> ssub * cmtt/m/it}{}
\newcommand{\texfamily}{\fontfamily{cmtex}\selectfont}
\DeclareFontShape{OT1}{cmtt}{bx}{n}
  {<5><6><7><8>cmtt8
   <9>cmbtt9
   <10><10.95><12><14.4><17.28><20.74><24.88>cmbtt10}{}
\DeclareFontShape{OT1}{cmtex}{bx}{n}
  {<-> ssub * cmtt/bx/n}{}
\newcommand{\tex}[1]{\text{\texfamily#1}}	% NEU

\newcommand{\Sp}{\hskip.33334em\relax}


\newcommand{\Conid}[1]{\mathit{#1}}
\newcommand{\Varid}[1]{\mathit{#1}}
\newcommand{\anonymous}{\kern0.06em \vbox{\hrule\@width.5em}}
\newcommand{\plus}{\mathbin{+\!\!\!+}}
\newcommand{\bind}{\mathbin{>\!\!\!>\mkern-6.7mu=}}
\newcommand{\rbind}{\mathbin{=\mkern-6.7mu<\!\!\!<}}% suggested by Neil Mitchell
\newcommand{\sequ}{\mathbin{>\!\!\!>}}
\renewcommand{\leq}{\leqslant}
\renewcommand{\geq}{\geqslant}
\usepackage{polytable}

%mathindent has to be defined
\@ifundefined{mathindent}%
  {\newdimen\mathindent\mathindent\leftmargini}%
  {}%

\def\resethooks{%
  \global\let\SaveRestoreHook\empty
  \global\let\ColumnHook\empty}
\newcommand*{\savecolumns}[1][default]%
  {\g@addto@macro\SaveRestoreHook{\savecolumns[#1]}}
\newcommand*{\restorecolumns}[1][default]%
  {\g@addto@macro\SaveRestoreHook{\restorecolumns[#1]}}
\newcommand*{\aligncolumn}[2]%
  {\g@addto@macro\ColumnHook{\column{#1}{#2}}}

\resethooks

\newcommand{\onelinecommentchars}{\quad-{}- }
\newcommand{\commentbeginchars}{\enskip\{-}
\newcommand{\commentendchars}{-\}\enskip}

\newcommand{\visiblecomments}{%
  \let\onelinecomment=\onelinecommentchars
  \let\commentbegin=\commentbeginchars
  \let\commentend=\commentendchars}

\newcommand{\invisiblecomments}{%
  \let\onelinecomment=\empty
  \let\commentbegin=\empty
  \let\commentend=\empty}

\visiblecomments

\newlength{\blanklineskip}
\setlength{\blanklineskip}{0.66084ex}

\newcommand{\hsindent}[1]{\quad}% default is fixed indentation
\let\hspre\empty
\let\hspost\empty
\newcommand{\NB}{\textbf{NB}}
\newcommand{\Todo}[1]{$\langle$\textbf{To do:}~#1$\rangle$}

\EndFmtInput
\makeatother
%
%
%
%
%
%
% This package provides two environments suitable to take the place
% of hscode, called "plainhscode" and "arrayhscode". 
%
% The plain environment surrounds each code block by vertical space,
% and it uses \abovedisplayskip and \belowdisplayskip to get spacing
% similar to formulas. Note that if these dimensions are changed,
% the spacing around displayed math formulas changes as well.
% All code is indented using \leftskip.
%
% Changed 19.08.2004 to reflect changes in colorcode. Should work with
% CodeGroup.sty.
%
\ReadOnlyOnce{polycode.fmt}%
\makeatletter

\newcommand{\hsnewpar}[1]%
  {{\parskip=0pt\parindent=0pt\par\vskip #1\noindent}}

% can be used, for instance, to redefine the code size, by setting the
% command to \small or something alike
\newcommand{\hscodestyle}{}

% The command \sethscode can be used to switch the code formatting
% behaviour by mapping the hscode environment in the subst directive
% to a new LaTeX environment.

\newcommand{\sethscode}[1]%
  {\expandafter\let\expandafter\hscode\csname #1\endcsname
   \expandafter\let\expandafter\endhscode\csname end#1\endcsname}

% "compatibility" mode restores the non-polycode.fmt layout.

\newenvironment{compathscode}%
  {\par\noindent
   \advance\leftskip\mathindent
   \hscodestyle
   \let\\=\@normalcr
   \let\hspre\(\let\hspost\)%
   \pboxed}%
  {\endpboxed\)%
   \par\noindent
   \ignorespacesafterend}

\newcommand{\compaths}{\sethscode{compathscode}}

% "plain" mode is the proposed default.
% It should now work with \centering.
% This required some changes. The old version
% is still available for reference as oldplainhscode.

\newenvironment{plainhscode}%
  {\hsnewpar\abovedisplayskip
   \advance\leftskip\mathindent
   \hscodestyle
   \let\hspre\(\let\hspost\)%
   \pboxed}%
  {\endpboxed%
   \hsnewpar\belowdisplayskip
   \ignorespacesafterend}

\newenvironment{oldplainhscode}%
  {\hsnewpar\abovedisplayskip
   \advance\leftskip\mathindent
   \hscodestyle
   \let\\=\@normalcr
   \(\pboxed}%
  {\endpboxed\)%
   \hsnewpar\belowdisplayskip
   \ignorespacesafterend}

% Here, we make plainhscode the default environment.

\newcommand{\plainhs}{\sethscode{plainhscode}}
\newcommand{\oldplainhs}{\sethscode{oldplainhscode}}
\plainhs

% The arrayhscode is like plain, but makes use of polytable's
% parray environment which disallows page breaks in code blocks.

\newenvironment{arrayhscode}%
  {\hsnewpar\abovedisplayskip
   \advance\leftskip\mathindent
   \hscodestyle
   \let\\=\@normalcr
   \(\parray}%
  {\endparray\)%
   \hsnewpar\belowdisplayskip
   \ignorespacesafterend}

\newcommand{\arrayhs}{\sethscode{arrayhscode}}

% The mathhscode environment also makes use of polytable's parray 
% environment. It is supposed to be used only inside math mode 
% (I used it to typeset the type rules in my thesis).

\newenvironment{mathhscode}%
  {\parray}{\endparray}

\newcommand{\mathhs}{\sethscode{mathhscode}}

% texths is similar to mathhs, but works in text mode.

\newenvironment{texthscode}%
  {\(\parray}{\endparray\)}

\newcommand{\texths}{\sethscode{texthscode}}

% The framed environment places code in a framed box.

\def\codeframewidth{\arrayrulewidth}
\RequirePackage{calc}

\newenvironment{framedhscode}%
  {\parskip=\abovedisplayskip\par\noindent
   \hscodestyle
   \arrayrulewidth=\codeframewidth
   \tabular{@{}|p{\linewidth-2\arraycolsep-2\arrayrulewidth-2pt}|@{}}%
   \hline\framedhslinecorrect\\{-1.5ex}%
   \let\endoflinesave=\\
   \let\\=\@normalcr
   \(\pboxed}%
  {\endpboxed\)%
   \framedhslinecorrect\endoflinesave{.5ex}\hline
   \endtabular
   \parskip=\belowdisplayskip\par\noindent
   \ignorespacesafterend}

\newcommand{\framedhslinecorrect}[2]%
  {#1[#2]}

\newcommand{\framedhs}{\sethscode{framedhscode}}

% The inlinehscode environment is an experimental environment
% that can be used to typeset displayed code inline.

\newenvironment{inlinehscode}%
  {\(\def\column##1##2{}%
   \let\>\undefined\let\<\undefined\let\\\undefined
   \newcommand\>[1][]{}\newcommand\<[1][]{}\newcommand\\[1][]{}%
   \def\fromto##1##2##3{##3}%
   \def\nextline{}}{\) }%

\newcommand{\inlinehs}{\sethscode{inlinehscode}}

% The joincode environment is a separate environment that
% can be used to surround and thereby connect multiple code
% blocks.

\newenvironment{joincode}%
  {\let\orighscode=\hscode
   \let\origendhscode=\endhscode
   \def\endhscode{\def\hscode{\endgroup\def\@currenvir{hscode}\\}\begingroup}
   %\let\SaveRestoreHook=\empty
   %\let\ColumnHook=\empty
   %\let\resethooks=\empty
   \orighscode\def\hscode{\endgroup\def\@currenvir{hscode}}}%
  {\origendhscode
   \global\let\hscode=\orighscode
   \global\let\endhscode=\origendhscode}%

\makeatother
\EndFmtInput
%
%
%
% First, let's redefine the forall, and the dot.
%
%
% This is made in such a way that after a forall, the next
% dot will be printed as a period, otherwise the formatting
% of `comp_` is used. By redefining `comp_`, as suitable
% composition operator can be chosen. Similarly, period_
% is used for the period.
%
\ReadOnlyOnce{forall.fmt}%
\makeatletter

% The HaskellResetHook is a list to which things can
% be added that reset the Haskell state to the beginning.
% This is to recover from states where the hacked intelligence
% is not sufficient.

\let\HaskellResetHook\empty
\newcommand*{\AtHaskellReset}[1]{%
  \g@addto@macro\HaskellResetHook{#1}}
\newcommand*{\HaskellReset}{\HaskellResetHook}

\global\let\hsforallread\empty

\newcommand\hsforall{\global\let\hsdot=\hsperiodonce}
\newcommand*\hsperiodonce[2]{#2\global\let\hsdot=\hscompose}
\newcommand*\hscompose[2]{#1}

\AtHaskellReset{\global\let\hsdot=\hscompose}

% In the beginning, we should reset Haskell once.
\HaskellReset

\makeatother
\EndFmtInput
% ----------------------
%% Stuff for TimCore


%% 'g' inferences





%% Intersection of effects

%% --------- Effects ------------


%% ----------------------


%% --------------------------
%%        Desugaring rules
%% --------------------------



% ----------------------------------------------------
% New desugaring, just does the wrapping part:
% ----------------------------------------------------
%   Input for wrapping


%%format Hide (e) = "{\mathcal H}\llbracket{}" e "\rrbracket{}"

%% 'si' is short for "solve or infer" or "unification variable or not" resp

%% 'md' is short for "mode" or "unification variable or not" resp

%% flipsi is a no-op for now


%% EXX is just EX, but with no argument




%% ----------- Applications --------------
%% %format applyx a b = a "@" b
%% %format applyx a b = a "\," b

%% Function application with no arguments

%% Function application with 1 argument

%% Function application with 1 argument, no parens

%% Function application with 1 argument, without parenthesis

%% Function application with many arguments

%% ----------- End of Applications --------------





%% Sequences e1; ...; en; v

%% xs etc stole syntax used in examples.

%% Old version: | for BNF, [] for choice
%%%format choice = "[ \! ]"
%%%format `choice` = "\mathop{[ \! ]}"
%% 1mm is about 0.3em with the current font

%% New version: tall | for BNF, fat | for choice
%%format vbar = "\mathop{\bigm|}"
%%format choice = "\boldsymbol{|}"
%%format `choice` = "\mathop{\boldsymbol{|}}"




%% %format ok = "\textbf{ok}"

%% %format `eqsemi` = "\hspace{-0.5mm}\fatsemi{}"
%% %format eqsemiKW = "\fatsemi"

%% Arrays and tuples

%% %format vmap (x) = "\llangle" x "\rrangle"






%%  OLD SYNTAX   %format split (x) (y) (z) = "\textbf{split}(" x ")\{" y "," z "\}"
%% %format defKW = "\textbf{def}"
% only used in comments, but still prettier this way:
%%%%format map = "\textbf{map}"

%% -------Functions ---------------


%% functionOC has open/closed subscript


%% -------End of f unctions ---------------





%% Unification
%%%format == = "\!=\!"


% format := = "\mapsto"
% %format squiggt = "\leadsto"
% %format squiggt = "\!-\!\!>\!\!\!"



































%% We combined S and C into one


%% ----------- Metatheory --------------


%% ----------------------------
%%         Verification
%% ----------------------------

%% assert = succeed, abbreviated to suc



%% Effects checking, like succeeds{e}





%% <fx> brackets

\subsection{Denotational semantics of Mini Verse}

%%%%%%%%%%%%%%%%%%%%%%%%%%%%%%%%%%%%%%%%%%%%%%%%%%%%%%%%%%%%%%%%%%%%%%%%%%%%%%%%%%%%%%%%%%%%%%%%%%%%%%%%%%%%%%%%%%%%%%%%%%%%%%%%%%%%%%%%%%%%%%%%%%%%%%%%%%%%%
%% \begin{figure}                                                                                                                                          %%
%% $$                                                                                                                                                      %%
%% \begin{array}{l}                                                                                                                                        %%
%%                                                                                                                                                         %%
%%                                                                                                                                                         %%
%% \textbf{Domains (sets)} \\                                                                                                                              %%
%% \begin{array}{rcll}                                                                                                                                     %%
%% \Value & = & \mathbb{Z} + (\semfcn) \\                                                                                                                  %%
%% \Env & = & |Ident| \rightarrow \Value \\                                                                                                                %%
%% W^{*} & = & \powerset{W} \\                                                                                                                             %%
%% Env^{*} & = & \powerset{Env} \\                                                                                                                         %%
%% \end{array}                                                                                                                                             %%
%% \\ \\                                                                                                                                                   %%
%% \begin{array}{r@@{\;\;\;}c@@{\;\;\;}l}                                                                                                                  %%
%%   \semap{E}{|e|} & : & \Env \rightarrow \mtype \\                                                                                                       %%
%%   \semapp{E}{|x|}{\rho} & = & \msing{\rho(x)} \\                                                                                                        %%
%%   \semapp{E}{|k|}{\rho} & = & \msing{k} \\                                                                                                              %%
%%   \semapp{E}{|o|}{\rho} & = & \msing{\semapp{O}{|o|}{}} \\                                                                                              %%
%%   \semapp{E}{|apply1N e_1 e_2|}{\rho} & = & %% \apply(\semapp{E}{|e_1|}{\rho}, \; \semapp{E}{|e_2|}{\rho}) \\                                           %%
%%     \{ r \, || \, f \inb \semapp{E}{|e_1|}{\rho}, \; a \inb \semapp{E}{|e_2|}{\rho}, \; r \inb \apply(f, a) \} \\                                       %%
%%   \semapp{E}{|e_1 = e_2|}{\rho}      & = & \semapp{E}{|e_1|}{\rho} \; \cap \; \semapp{E}{|e_2|}{\rho} \\                                                %%
%%   \semapp{E}{|e_1 ; e_2|}{\rho}      & = & %% \semapp{E}{|e_1|}{\rho} \; \msemi      \; \semapp{E}{|e_2|}{\rho} \\                                      %%
%%     \{ y \, || \, x \inb \semapp{E}{|e_1|}{\rho}, \; y \inb \semapp{E}{|e_2|}{\rho} \} \\                                                               %%
%%   \semapp{E}{|defKW x|}{\rho} & = & \msing{|emptytup|} \\                                                                                               %%
%%   \semapp{E}{|fail|}{\rho} & = & \mempty \\                                                                                                             %%
%%   \semapp{E}{|arr (e_0, xdots, e_n)|}{\rho} & = &                                                                                                       %%
%%     \{ \mkfcn{ \{ (i,x_i) \vb i \in 0 \ldots n \} } \, || \, x_0 \inb \semapp{E}{|e_1|}{\rho}, \,\mydots,\, x_n \inb \semapp{E}{|e_n|}{\rho} \} \\[1mm] %%
%%   \semapp{E}{|vif e_1 e_2 e_3|}{\rho} & = &                                                                                                             %%
%%   \left\{                                                                                                                                               %%
%%   \begin{array}{ll}                                                                                                                                     %%
%%     \semapp{D}{|e_3|}{\rho} & \text{if}\; \bar{\rho} = \mempty \\                                                                                       %%
%%     \bigcup_{\rho' \dot{\in} \bar{\rho}} \semapp{D}{|e_2|}{\rho'} & \text{otherwise} \\                                                                 %%
%%   \end{array}                                                                                                                                           %%
%%   \right. \\                                                                                                                                            %%
%%   & &  |where|\; \bar{\rho} = \semappp{C}{e_1}{\rho} \\[1mm]                                                                                            %%
%%   \textit{succeeds function:} \\                                                                                                                        %%
%%   \semapp{E}{|lamOC q x e_1 e_2|}{\rho} & = & \semclose{q}{\{ f \in \semfcn \;                                                                          %%
%%     \begin{array}[t]{@@{}l@@{\;}l}                                                                                                                      %%
%%       || & \forall x \in \Value.  \forall \rho' \in \semapp{X}{e_1}{\rho}. \\                                                                           %%
%%         & \hspace{1em} (\semapp{E}{e_1}{\rho'[x \fmapsto w]} \neq \mempty) \\                                                                           %%
%%         & \hspace{1em} \implies x \in \dom(f) \land f(x) \in \semapp{D}{e_2}{\rho'}} \}  \\                                                             %%
%%     \end{array} \\                                                                                                                                      %%
%%   \textit{decides function:} \\                                                                                                                         %%
%%   \semapp{E}{|lamOC q x e_1 e_2|}{\rho} & = & \semclose{q}{\{ f \in \semfcn \;                                                                          %%
%%     \begin{array}[t]{@@{}l@@{\;}l}                                                                                                                      %%
%%       || & \forall x \in \Value.  \forall \rho' \in \semapp{X}{e_1}{\rho}. \\                                                                           %%
%%         & \hspace{1em} (\semapp{E}{e_1}{\rho'[x \fmapsto w]} \neq \mempty \land \semapp{D}{e_2}{\rho'} \neq \mempty) \\                                 %%
%%         & \hspace{1em} \implies x \in \dom(f) \land f(x) \in \semapp{D}{e_2}{\rho'}} \}  \\                                                             %%
%%     \end{array} \\                                                                                                                                      %%
%%                                                                                                                                                         %%
%%                                                                                                                                                         %%
%% \end{array}                                                                                                                                             %%
%% \\                                                                                                                                                      %%
%% \begin{array}{r@@{\;\;\;}c@@{\;\;\;}ll}                                                                                                                 %%
%% \semap{D}{|e|} & : & \Env \rightarrow \mtype \\                                                                                                         %%
%% \semapp{D}{|e|}{\rho} & = & \bigcup_{\rho' \dot{\in} \semapp{X}{|e|}{\rho}} \semapp{E}{|e|}\rho' \\[2mm]                                                %%
%% \semap{X}{e} & : & \Env \rightarrow \powerset{\Env} \\                                                                                                  %%
%% \semapp{X}{e}{\rho} & = & \{ \extendrho{\rho}{x_1 \fmapsto w_1, \mydots, x_n \fmapsto w_n }                                                             %%
%%                                   \vb   w_1 \inb \Value, \mydots, w_n \inb \Value \} \\                                                                 %%
%%   & & |where [x_1,xdots,x_n]| = \semap{I}{|e|} \\[1mm]                                                                                                  %%
%%                                                                                                                                                         %%
%% \semap{I}{|e|} & : & [\mathit{Ident}] \\                                                                                                                %%
%% \semap{I}{|e|} & = & \text{the outer $\exists$-bound variables in |e|} \\[1mm]                                                                          %%
%%                                                                                                                                                         %%
%%   \semapp{O}{|op|}{} & : & \semfcn \\                                                                                                                   %%
%%   \semapp{O}{|gt|}{} & = & \fcn \{ (\tuple{x,y},x) \, || \, x \inb \mathbb{Z}, y \inb \mathbb{Z}, x > y \} \\                                           %%
%%   \semapp{O}{|int|}{} & = & \fcn \{ (x,x) \, || \, x \inb \mathbb{Z} \} \} \\                                                                           %%
%%   \semapp{O}{|add|}{} & = & \fcn \{ (\tuple{x,y},x+y) \, || \, x \inb \mathbb{Z}, y \inb \mathbb{Z} \} \\                                               %%
%% \hline \\[-2mm]                                                                                                                                         %%
%% \apply & : & (\Value \times \Value) \rightarrow \mtype \\                                                                                               %%
%% \apply( f, w ) & = & \msing{f(w)} & f \in \semfcn, w \in \dom(f) \\                                                                                     %%
%% \apply( f, w ) & = & \mempty & \text{otherwise} \\[1mm]                                                                                                 %%
%%                                                                                                                                                         %%
%% \semclose{|open|}{F} & = & F \\                                                                                                                         %%
%% \semclose{|closed|}{F} & = & \{ f \inb F \vb \forall f' \in F. \dom{f} \subseteq \dom{f'} \} \\[1mm]                                                    %%
%%                                                                                                                                                         %%
%% \fcn & : & (\tuple{\Value,\Value})^{*} \rightarrow (\semfcn) \\                                                                                         %%
%% \fcn & & \multicolumn{2}{l}{\text{indicates a function made from a set of pairs} } \\                                                                   %%
%%                                                                                                                                                         %%
%%                                                                                                                                                         %%
%% \end{array}                                                                                                                                             %%
%% \end{array}                                                                                                                                             %%
%% $$                                                                                                                                                      %%
%% \caption{Plan C: Idealised denotational semantics of Mini Verse (without choice)}                                                                       %%
%% \end{figure}                                                                                                                                            %%
%%%%%%%%%%%%%%%%%%%%%%%%%%%%%%%%%%%%%%%%%%%%%%%%%%%%%%%%%%%%%%%%%%%%%%%%%%%%%%%%%%%%%%%%%%%%%%%%%%%%%%%%%%%%%%%%%%%%%%%%%%%%%%%%%%%%%%%%%%%%%%%%%%%%%%%%%%%%%

% ------------- Denotional semantics of Mini Verse with choice -------------------

\subsection{Denotional semantics of Mini Verse with choice}

\newcommand{\applyo}{\textbf{applyo}}
\newcommand{\applyf}{\textbf{applyf}}
\newcommand{\applys}{\textbf{applys}}
\newcommand{\squash}{\textbf{squash}}
\newcommand{\Ws}{\textit{Ws}}
\newcommand{\mkall}{\textbf{all}}
\newcommand{\mkarr}{\textbf{arr}}
\newcommand{\combine}{\textbf{comb}}
\newcommand{\trunc}{\textbf{trunc}}
\newcommand{\once}{\textbf{once}}
\newcommand{\sequence}{\textbf{seq}}
\newcommand{\maplist}{\textbf{map}}
%% \newcommand{\maplist}{\textbf{map}_{[]}}
\newcommand{\mapset}{\textbf{map}_{\{\}}}
\newcommand{\cross}{\textbf{cross}}
\newcommand{\dist}{\textbf{dist}}
\newcommand{\group}{\textbf{group}}
\newcommand{\myset}[1]{\text{Set}(#1)}
%\newcommand{\fold}{\textbf{fold}}
%\newcommand{\filter}{\textbf{filter}}
\newcommand{\myenv}[1]{E\{ #1 \}}
\newcommand{\mybind}[2]{#1 \mapsto #2}
\newcommand{\succeeds}{\textbf{succeeds}}
\newcommand{\seqcompat}{\textbf{compat}}
\newcommand{\setcompat}{\textbf{comp}}

\begin{figure}
$$
\begin{array}{l}

\textbf{Domains (sets)} \\
\begin{array}{rcll}
\Value & = & \mathbb{Z} + [\semfcn] \\
\Env & = & \ensuremath{\Conid{Ident}} \rightarrow \Value \\
\Ws & = & \powerset{W} \\
W^{*} & = & [\Ws] \\
Env^{*} & = & \myset{Env} \\
\end{array}
\\ \\
\begin{array}{r@{\;\;\;}c@{\;\;\;}l}
  \semap{E}{\ensuremath{\Varid{e}}} & : & \Env \rightarrow \mtype \\
  \semapp{E}{\ensuremath{\Varid{x}}}{\rho} & = & [\msing{\rho(x)}] \\
  \semapp{E}{\ensuremath{\Varid{k}}}{\rho} & = & [\msing{k}] \\
  \semapp{E}{\ensuremath{\Varid{o}}}{\rho} & = & [\msing{\semapp{O}{\ensuremath{\Varid{o}}}{}}] \\
  \semapp{E}{\ensuremath{\exists{}\;\Varid{x}}}{\rho} & = & [\msing{\ensuremath{\langle\rangle}}] \\
  \semapp{E}{\ensuremath{\textbf{fail}}}{\rho} & = & [] \\
  \semapp{E}{\ensuremath{\Varid{e}_{\mathrm{1}}\hspace{0.19em}\mathop{\rule[-0.1em]{0.15em}{0.75em}}\hspace{0.2em}\Varid{e}_{\mathrm{2}}}}{\rho} & = & \semapp{D}{\ensuremath{\Varid{e}_{\mathrm{1}}}}{\rho} \ensuremath{\bowtie} \semapp{D}{\ensuremath{\Varid{e}_{\mathrm{1}}}}{\rho} \\
  \semapp{E}{\ensuremath{\Varid{e}_{\mathrm{1}} [\!\Varid{e}_{\mathrm{2}}]}}{\rho} & = &
         [ r \, | \,
           \ensuremath{\Varid{fs}\leftarrow } \semapp{E}{\ensuremath{\Varid{e}_{\mathrm{1}}}}{\rho}, \;
            as\ensuremath{\leftarrow } \semapp{E}{\ensuremath{\Varid{e}_{\mathrm{2}}}}{\rho}, \;
           \ensuremath{\Varid{r}} \inb \applys(fs, as) \}
         ]
           \\

  \semapp{E}{\ensuremath{\Varid{e}_{\mathrm{1}}\mathrel{=}\Varid{e}_{\mathrm{2}}}}{\rho}      & = &
      [ s_1 \cap s_2\, |\, 
           \ensuremath{\Varid{s}_{\mathrm{1}}\leftarrow } \semapp{E}{\ensuremath{\Varid{e}_{\mathrm{1}}}}{\rho}, \; \ensuremath{\Varid{s}_{\mathrm{2}}\leftarrow } \semapp{E}{\ensuremath{\Varid{e}_{\mathrm{2}}}}{\rho} ] \\

  \semapp{E}{\ensuremath{\Varid{e}_{\mathrm{1}};\,\Varid{e}_{\mathrm{2}}}}{\rho}      & = &
    [ \{ x_2 \, | \, x_1 \inb s_1, \; x_2 \inb s_2 \}\, |\,
      \ensuremath{\Varid{s}_{\mathrm{1}}\leftarrow } \semapp{E}{\ensuremath{\Varid{e}_{\mathrm{1}}}}{\rho}, \; \ensuremath{\Varid{s}_{\mathrm{2}}\leftarrow } \semapp{E}{\ensuremath{\Varid{e}_{\mathrm{2}}}}{\rho} ] \\

  \semapp{E}{\ensuremath{\langle\Varid{e}_{\mathrm{0}},\mydots,\Varid{e}_{\Varid{n}}\rangle}}{\rho} & = &
    [ \{ \mkarr(x_0,\mydots,x_n) \, |\, x_0 \inb s_0,\; \mydots,\; x_n \inb s_n \}\, \\
      & & |\, s_0 \ensuremath{\leftarrow } \semapp{E}{\ensuremath{\Varid{e}_{\mathrm{0}}}}{\rho}, \,\mydots,\; s_n \ensuremath{\leftarrow } \semapp{E}{\ensuremath{\Varid{e}_{\Varid{n}}}}{\rho} ] \\

  \semapp{E}{\ensuremath{\textbf{all}\{\Varid{e}\}}}{\rho} & = & [ \mkall(\semapp{E}{e}{\rho}) ] \\[1mm]

  \semapp{E}{\ensuremath{\textbf{if}(\Varid{e}_{\mathrm{1}})\{\Varid{e}_{\mathrm{2}}\}\textbf{else}\{\Varid{e}_{\mathrm{3}}\}}}{\rho} & = &
  \left\{
  \begin{array}{ll}
    \semapp{D}{\ensuremath{\Varid{e}_{\mathrm{3}}}}{\rho} & \text{if}\; \bar{\rho} = \mempty \\
    \squash(\hat{\bigcup}\{ \semapp{D}{e_2}{\rho'} \, | \, \rho' \inb \bar{\rho} \} & \text{otherwise} \\
  \end{array}
  \right. \\
  & &  \ensuremath{\mathbf{where}}\; \bar{\rho} = \semapp{R}{e_1}{\rho}
  \\[1mm]

%  \textit{succeeds function:} \\
%    \semapp{E}{|lamOC q x e_1 e_2|}{\rho} & = &
%    \left\{
%    \begin{array}{ll}
%      [\semopen{q}(\combine(|fs|))] & \text{if } \succeeds(|fs|) \\
%      \mbox{}[\mempty] & \text{otherwise}\\
%    \end{array}
%    \right. \\
%    & & |where|
%      \begin{array}[t]{r@@{\;}c@@{\;}l@@{\;}l}
%        |fs| & = &
%        \{ [ \dist(v, s) \,||\, s \inb r ] \\
%        & & ||\, v \inb \Value, |let r| = \semapp{F}{|lamOC q x e_1 e_2|}{\rho}\;v, r \neq [] \} \\
%      \end{array}
    \semapp{E}{\ensuremath{\lambda_{\Varid{q}}\Varid{x}.(\Varid{e}_{\mathrm{1}})\{\Varid{e}_{\mathrm{2}}\}}}{\rho} & = & [ \semopen{q}{(fs)} ] \\
    & & \ensuremath{\mathbf{where}}
      \begin{array}[t]{r@{\;}c@{\;}l@{\;}l}
        \ensuremath{\Varid{fs}} & = & \{ f \,|\, f \inb [\semfcn], \\
        & & \forall v \inb \Value .\\
        & & \seqcompat(\applyo(f,v), \semappp{F}{\ensuremath{\lambda_{\Varid{q}}\Varid{x}.(\Varid{e}_{\mathrm{1}})\{\Varid{e}_{\mathrm{2}}\}}}{\rho}{v}) \}
      \end{array}
    
  \\
  \\
\semap{F}{e} & : & \Env \rightarrow \Value \rightarrow [\myset{\Value}] \\
\semapp{F}{\ensuremath{\lambda_{\Varid{q}}\Varid{x}.(\Varid{e}_{\mathrm{1}})\{\Varid{e}_{\mathrm{2}}\}}}{\rho}\;v & = & \trunc(r, \ensuremath{\Varid{ws}}) \\
& & \ensuremath{\mathbf{where}}
    \begin{array}[t]{r@{\;}c@{\;}l@{\;}l}
      \ensuremath{\Varid{ws}} & = & \hat{\bigcup} \{ & [ \ensuremath{\textbf{if}\;\Varid{s}\mathrel{=}}\mempty\,\ensuremath{\textbf{then}}\,\mempty\,\ensuremath{\textbf{else}}\, \ensuremath{\Varid{w}}\,|\, s \inb ss ] \\
      & & & |\, \rho' \inb \semapp{X}{\ensuremath{\Varid{e}_{\mathrm{1}}}}{(\rho[x \fmapsto v])}, \\
      & & & \ensuremath{\mathbf{let}\;\Varid{ss}} = \semapp{E}{\ensuremath{\Varid{e}_{\mathrm{1}}}}{\rho'},\, \squash(\ensuremath{\Varid{ss}}) \neq [], \\
      & & & \ensuremath{\mathbf{let}\;\Varid{w}} = \once(\semapp{D}{e_2}{\rho'}) \, \} \\
      \ensuremath{\Varid{r}} & = & \multicolumn{2}{l} { \bigcap [ w \, | \, w \inb ws, ws \neq \mempty ] }
    \end{array} \\


%\end{array}
\\ 
%\begin{array}{r@@{\;\;\;}c@@{\;\;\;}ll}
\semap{D}{\ensuremath{\Varid{e}}} & : & \Env \rightarrow \mtype \\
%\semapp{D}{|e|}{\rho} & = & \hat{\bigcup}_{\rho' \dot{\in} \semapp{X}{|e|}{\rho}} \semapp{E}{|e|}\rho' \\[2mm]
\semapp{D}{\ensuremath{\Varid{e}}}{\rho} & = & \hat{\bigcup}\{ \semapp{E}{e}{\rho'} \, | \, \rho' \inb \semapp{X}{e}{\rho} \} \\[2mm]
\semap{R}{\ensuremath{\Varid{e}}}{} & : & \Env \rightarrow \myset{\Env} \\
\semapp{R}{\ensuremath{\Varid{e}}}{\rho} & = & \{ \rho' \, | \, \rho' \inb \semapp{X}{\ensuremath{\Varid{e}}}{\rho}, \squash(\semapp{E}{\ensuremath{\Varid{e}}}{\rho'}) \neq [] \} \\[2mm]
%% \semap{C}{|e|} & : & \Env -> \rightarrow [Env^{*}] \\
%% \semapp{C}{|e|}{\rho} & = & xxx \\
\semap{X}{e} & : & \Env \rightarrow \myset{\Env} \\
\semapp{X}{e}{\rho} & = & \{ \extendrho{\rho}{x_1 \fmapsto w_1, \mydots, x_n \fmapsto w_n }
                                  \vb   w_1 \inb \Value, \mydots, w_n \inb \Value \} \\
  & & \ensuremath{\mathbf{where}\;[\mskip1.5mu \Varid{x}_{\mathrm{1}},\mydots,\Varid{x}_{\Varid{n}}\mskip1.5mu]} = \semap{I}{\ensuremath{\Varid{e}}} \\[1mm]

\semap{I}{\ensuremath{\Varid{e}}} & : & [\mathit{Ident}] \\
\semap{I}{\ensuremath{\Varid{e}}} & = & \text{the outer $\exists$-bound variables in \ensuremath{\Varid{e}}} \\[1mm]
\end{array}
\\

\end{array}
$$
\caption{Plan C: Idealised denotational semantics of Mini Verse with choice}
\end{figure}


\begin{figure}
$$
\begin{array}{l}

\textbf{Auxiliary functions}\\
\begin{array}{l@{\;}c@{\;}ll}
\apply & : & (\Value \times \Value) \rightarrow \mtype \\
\apply( f, w ) & = & \msing{f(w)} & f \in \semfcn, w \in \dom(f) \\
\apply( f, w ) & = & \mempty & \text{otherwise} \\[1mm]

  \applys & : & (\Ws, \Ws) \rightarrow [\Ws] \\
  \applys(\ensuremath{\Varid{fss}}, as) & = & \hat{\bigcup} \{ \applyf(\ensuremath{\Varid{fs}}, as) \, | \, \ensuremath{\Varid{fs}} \inb \ensuremath{\Varid{fss}} \} \\
  \applyf & : & ([\semfcn], \Ws) \rightarrow [\Ws] \\
  \applyf(\ensuremath{\Varid{fs}}, as) & = & \bar{\bigcup} \{ \applyo(\ensuremath{\Varid{fs}}, a) \, | \, a \inb as \} \\
  \applyo & : & ([\semfcn], \Value) \rightarrow [\Ws] \\
  \applyo(\ensuremath{\Varid{fs}},a) & = & [ \apply(\ensuremath{\Varid{f},\Varid{a}}) \, | \, \ensuremath{\Varid{f}} \inb \ensuremath{\Varid{fs}} ] \\[2mm]
  \squash & : & \mtype \rightarrow \mtype \\
  \squash(xs) & = & [ x \, | \, x \inb xs, x \neq \{\} ] & \text{remove empty sets from the sequence} \\
  \mkall([s_1,\mydots,s_n]) & = & \{ \mkarr(x_1,\mydots,x_n) \, | \, x_1 \inb s_1, \, \mydots \, x_n \inb s_n \} \\
  \mkarr(x_0,\mydots,x_n) & = & [ \{0 \mapsto x_0\}, \mydots, \{n \mapsto x_n\} ]
                              & \text{make an array-function} \\
  \once([x]) & = & x \\
  \once(\_) & = & \mempty & \text{or WRONG} \\
  \trunc(x,[\mempty, \mydots, s, \mydots]) & = & [\mempty, \mydots, x] & s \neq \mempty, \text{replace first non-empty} \\
\\

\semopen{q} & : & [\semfcn] \rightarrow [\semfcn] \\
\semopen{\ensuremath{\textbf{o}}}{(G)} & = & [ \{ f \in \semfcn\, |\, \exists g \inb \bigcup G\, .\, g \subseteq f  \} ] \\
\semopen{\ensuremath{\textbf{c}}}{(G)} & = & G \\[1mm]

% \succeeds & : & \myset{[\myset{\Value}]} \rightarrow Bool \\
% \succeeds(s) & = & \multicolumn{2}{l}{\text{if all sequences in $s$ have a non-empty element}} \\[1mm]

\bigcap[] & = & \mempty \\
\bigcap[x_1,x_2,\mydots,x_n] & = & x_1 \cap x_2 \cap \mydots \cap x_n \\[1mm]

\seqcompat & : & (\mtype \times \mtype) \rightarrow \ensuremath{\Conid{Bool}} \\
\seqcompat([s_1,\mydots,s_n], [t_1,\mydots,t_m]) & = & n = m \land \setcompat(s_1,t_1) \land \mydots \land \setcompat(s_n,t_n) \\
\setcompat & : & (\myset{\Value} \times \myset{\Value}) \rightarrow Bool \\
\setcompat(\{\},\{\}) & = & \ensuremath{\textbf{true}} \\
\setcompat(\{\},\{\mydots\}) & = & \ensuremath{\textbf{false}} \\
\setcompat(\{\mydots\},\{\}) & = & \ensuremath{\textbf{false}} \\
\setcompat(s, t) & = & s \subseteq t \\[1mm]

\end{array} \\ \\

\textbf{Notation}\\
\begin{array}{lcll}
  \ensuremath{[\mskip1.5mu \Varid{e}_{\mathrm{1}},\mydots,\Varid{e}_{\Varid{n}}\mskip1.5mu]} & : & \mtype & \text{sequence of items} \\[2mm]

  \ensuremath{[\mskip1.5mu \Varid{e}\mid \Varid{x}_{\mathrm{1}}\leftarrow \Varid{e}_{\mathrm{1}},\Varid{x}_{\mathrm{2}}\leftarrow \Varid{e}_{\mathrm{2}},\mydots\mskip1.5mu]} & : & \mtype & \text{sequence comprehension (order preserving)}, \\[2mm]

  \ensuremath{\bowtie} & : & \mtype \rightarrow \mtype \rightarrow \mtype & \text{sequence concatenation} \\
  \ensuremath{[\mskip1.5mu \Varid{x}_{\mathrm{1}},\mydots,\Varid{x}_{\Varid{n}}\mskip1.5mu]\bowtie[\mskip1.5mu \Varid{y}_{\mathrm{1}},\mydots,\Varid{y}_{\Varid{n}}\mskip1.5mu]} & = & \ensuremath{[\mskip1.5mu \Varid{x}_{\mathrm{1}},\mydots,\Varid{x}_{\Varid{n}},\Varid{y}_{\mathrm{1}},\mydots,\Varid{y}_{\Varid{n}}\mskip1.5mu]} \\[2mm]

  \hat{\cup} & : & \mtype \rightarrow \mtype \rightarrow \mtype & \text{elementwise union of sequences,} \\
  \ensuremath{[\mskip1.5mu \Varid{x}_{\mathrm{1}},\mydots,\Varid{x}_{\Varid{n}}\mskip1.5mu]} \hat{\cup} \ensuremath{[\mskip1.5mu \Varid{y}_{\mathrm{1}},\mydots,\Varid{y}_{\Varid{n}}\mskip1.5mu]} & = & \ensuremath{[\mskip1.5mu \Varid{x}_{\mathrm{1}}\mathop{\cup}\Varid{y}_{\mathrm{1}},\mydots,\Varid{x}_{\Varid{n}}\mathop{\cup}\Varid{y}_{\Varid{n}}\mskip1.5mu]} & \text{right-padding shorter sequence with \mempty} \\[2mm]

  \hat{\bigcup} & : & \powerset{\mtype} \rightarrow \mtype \\
  \hat{\bigcup} & = & \cdots & \hat{\cup} \text{ extended to sets} \\[2mm]

%  |re| & : & \mtype \rightarrow \mtype & \text{remove all empty sets} \\
%  |re(s)| & = & ... \\
\end{array}

\end{array}
$$
\caption{Plan C: Idealised denotational semantics of Mini Verse with choice, notation}
\end{figure}

\newcommand{\newfcn}{\Value \hookrightarrow \Value}
\newcommand{\mknewfcn}[1]{[ #1 ]}
\newcommand{\perm}[1]{\textbf{perm}(#1)}

%%%%%%%%%%%%%%%%%%%%%%%%%%%%%%%%%%%%%%%%%%%%%%%%%%%%%%%%%%%%%%%%%%%%%%%%%%%%%%%%%%%%%%%%%%%%%%%%%%%%%%%%%%%%%%%%%%%%%%%%%%%%%%%%%%%
%% \begin{figure}                                                                                                                %%
%% $$                                                                                                                            %%
%% \begin{array}{l}                                                                                                              %%
%%                                                                                                                               %%
%% \textbf{Domains (sets)} \\                                                                                                    %%
%% \begin{array}{rcll}                                                                                                           %%
%% \Value & = & \mathbb{Z} + (\newfcn) \\                                                                                        %%
%% \newfcn & = & [\Value \times \Value] & \text{an orderered set of pairs} \\                                                    %%
%% \Env & = & |Ident| \rightarrow \Value \\                                                                                      %%
%% \Ws & = & \powerset{\Value} \\                                                                                                %%
%% W^{*} & = & [\Ws] \\                                                                                                          %%
%% Env^{*} & = & \powerset{Env} \\                                                                                               %%
%% \end{array}                                                                                                                   %%
%% \\ \\                                                                                                                         %%
%% \begin{array}{r@@{\;\;\;}c@@{\;\;\;}ll}                                                                                       %%
%%   \semap{E}{|e|} & : & \Env \rightarrow \mtype \\                                                                             %%
%%   \semapp{E}{|x|}{\rho} & = & [\msing{\rho(x)}] \\                                                                            %%
%%   \semapp{E}{|k|}{\rho} & = & [\msing{k}] \\                                                                                  %%
%%   \semapp{E}{|o|}{\rho} & = & [\perm{\semap{O}{|o|}}] \\                                                                      %%
%%   \semapp{E}{|defKW x|}{\rho} & = & [\msing{|emptytup|}] \\                                                                   %%
%%   \semapp{E}{|fail|}{\rho} & = & [\{\}] \\                                                                                    %%
%%   \semapp{E}{|e_1 `choice` e_2|}{\rho} & = & \semapp{D}{|e_1|}{\rho} |++| \semapp{D}{|e_1|}{\rho} \\                          %%
%%   \semapp{E}{|apply1N e_1 e_2|}{\rho} & = &                                                                                   %%
%%          [ r \, || \,                                                                                                         %%
%%            |fs <-| \semapp{E}{|e_1|}{\rho}, \;                                                                                %%
%%            |as <-| \semapp{E}{|e_2|}{\rho}, \;                                                                                %%
%%            |r| \inb \applys(|fs|, as) \}                                                                                      %%
%%          ]                                                                                                                    %%
%%            \\                                                                                                                 %%
%%                                                                                                                               %%
%%   \semapp{E}{|e_1 = e_2|}{\rho}      & = &                                                                                    %%
%%       [ s_1 \cap s_2\, ||\,                                                                                                   %%
%%            |s_1 <-| \semapp{E}{|e_1|}{\rho}, \; |s_2 <-| \semapp{E}{|e_2|}{\rho} ] \\                                         %%
%%                                                                                                                               %%
%%   \semapp{E}{|e_1 ; e_2|}{\rho}      & = &                                                                                    %%
%%     [ \{ x_2 \, || \, x_1 \inb s_1, \; x_2 \inb s_2 \}\, ||\,                                                                 %%
%%       |s_1 <-| \semapp{E}{|e_1|}{\rho}, \; |s_2 <-| \semapp{E}{|e_2|}{\rho} ] \\                                              %%
%%                                                                                                                               %%
%%   \semapp{E}{|arr (e_0, xdots, e_n)|}{\rho} & = &                                                                             %%
%%     [ \{ \mknewfcn{ (0,x_0),(1,x_1),\mydots,(n,x_n) } \, ||\, x_0 \inb s_0,\; \mydots,\; x_n \inb s_n \}\, \\                 %%
%%       & & ||\, s_0 |<-| \semapp{E}{|e_0|}{\rho}, \,\mydots,\; s_n |<-| \semapp{E}{|e_n|}{\rho} ] \\[1mm]                      %%
%%                                                                                                                               %%
%%   \semapp{E}{|vif e_1 e_2 e_3|}{\rho} & = &                                                                                   %%
%%   \left\{                                                                                                                     %%
%%   \begin{array}{ll}                                                                                                           %%
%%     \semapp{D}{|e_3|}{\rho} & \text{if}\; \bar{\rho} = \mempty \\                                                             %%
%%     \squash(\hat{\bigcup}\{ \semapp{D}{e_2}{\rho'} \, || \, \rho' \inb \bar{\rho} \} & \text{otherwise} \\                    %%
%%   \end{array}                                                                                                                 %%
%%   \right. \\                                                                                                                  %%
%%   & &  |where|\; \bar{\rho} = \semapp{R}{e_1}{\rho}                                                                           %%
%%   \\[1mm]                                                                                                                     %%
%%                                                                                                                               %%
%% %  \textit{succeeds function:} \\                                                                                             %%
%%   \semapp{E}{|lamOC q x e_1 e_2|}{\rho} & = & XXX \\                                                                          %%
%%   & & |where|                                                                                                                 %%
%%     \begin{array}[t]{r@@{\;}c@@{\;}l}                                                                                         %%
%%       |xxx| & = & |yyy| \\                                                                                                    %%
%%     \end{array} \\                                                                                                            %%
%%                                                                                                                               %%
%%                                                                                                                               %%
%% %    \semclose{q}                                                                                                             %%
%% %      {\{ f\,                                                                                                                %%
%% %    \begin{array}[t]{@@{}l@@{\;}l}                                                                                           %%
%% %      || & \forall x \inb \Value,\,\\                                                                                        %%
%% %         & \forall \rho' \in \semapp{X}{e_1}{\rho}, \\                                                                       %%
%% %         & f \inb \semfcn, \\                                                                                                %%
%% %        & \hspace{1em} (\semapp{E}{e_1}{\rho'[x \fmapsto w]} \neq \mempty) \\                                                %%
%% %        & \hspace{1em} \implies x \in \dom(f) \land f(x) \in \semapp{D}{e_2}{\rho'}} \}  \\                                  %%
%% %    \end{array} \\                                                                                                           %%
%% %  \textit{decides function:} \\                                                                                              %%
%% %  \semapp{E}{|lamOC q x e_1 e_2|}{\rho} & = & \semclose{q}{\{ f \in \semfcn \;                                               %%
%% %    \begin{array}[t]{@@{}l@@{\;}l}                                                                                           %%
%% %      || & \forall x \in \Value.  \forall \rho' \in \semapp{X}{e_1}{\rho}. \\                                                %%
%% %        & \hspace{1em} (\semapp{E}{e_1}{\rho'[x \fmapsto w]} \neq \mempty \land \semapp{D}{e_2}{\rho'} \neq \mempty) \\      %%
%% %        & \hspace{1em} \implies x \in \dom(f) \land f(x) \in \semapp{D}{e_2}{\rho'}} \}  \\                                  %%
%% %    \end{array} \\                                                                                                           %%
%%                                                                                                                               %%
%%                                                                                                                               %%
%% %\end{array}                                                                                                                  %%
%% \\ \\                                                                                                                         %%
%% %\begin{array}{r@@{\;\;\;}c@@{\;\;\;}ll}                                                                                      %%
%% \semap{D}{|e|} & : & \Env \rightarrow \mtype \\                                                                               %%
%% %\semapp{D}{|e|}{\rho} & = & \hat{\bigcup}_{\rho' \dot{\in} \semapp{X}{|e|}{\rho}} \semapp{E}{|e|}\rho' \\[2mm]               %%
%% \semapp{D}{|e|}{\rho} & = & \hat{\bigcup}\{ \semapp{E}{e}{\rho'} \, || \, \rho' \inb \semapp{X}{e}{\rho} \} \\[2mm]           %%
%% \semap{R}{|e|}{} & : & \Env \rightarrow \powerset{\Env} \\                                                                    %%
%% \semapp{R}{|e|}{\rho} & = & \{ \rho' \, || \, \rho' \inb \semapp{X}{|e|}{\rho}, \squash(\semapp{E}{|e|}{\rho'}) \neq [] \} \\ %%
%% %% \semap{C}{|e|} & : & \Env -> \rightarrow [Env^{*}] \\                                                                      %%
%% %% \semapp{C}{|e|}{\rho} & = & xxx \\                                                                                         %%
%% \end{array}                                                                                                                   %%
%% \\                                                                                                                            %%
%% \\                                                                                                                            %%
%% \begin{array}{rcll}                                                                                                           %%
%%   \applys & : & (\Ws, \Ws) \rightarrow [\Ws] \\                                                                               %%
%%   \applys(|fs|, as) & = & \hat{\bigcup} \{ \applyf|(f, a)| \, || \, |f| \inb |fs|, \, |f| \in (\newfcn), \, a \inb as \} \\   %%
%%   \applyf & : & (\newfcn, \Value) \rightarrow [\Ws] \\                                                                        %%
%%   \applyf(|f|, a) & = & [ \{ y \, || \, x=a \} \,  || \, (x, y) \inb f ] \\                                                   %%
%%   & & \text{sequence of sets, all empty except at most one } \\                                                               %%
%%   \perm{} & : & (\semfcn) \rightarrow \powerset{\newfcn} \\                                                                   %%
%%   \perm{f} & = & \text{take an ordinary function and produce all permutations of argument order} \\                           %%
%% \end{array} \\                                                                                                                %%
%% \lennart{I think these are correct for Tim's new functions.  Lambdas still to be done.}                                       %%
%% \\                                                                                                                            %%
%% \end{array}                                                                                                                   %%
%% $$                                                                                                                            %%
%% \caption{Plan C: Idealised denotational semantics of Mini Verse with choice, new functions}                                   %%
%% \end{figure}                                                                                                                  %%
%%%%%%%%%%%%%%%%%%%%%%%%%%%%%%%%%%%%%%%%%%%%%%%%%%%%%%%%%%%%%%%%%%%%%%%%%%%%%%%%%%%%%%%%%%%%%%%%%%%%%%%%%%%%%%%%%%%%%%%%%%%%%%%%%%%

The semantic function $\semap{F}{}$ gives meaning to a function applied to a single value, \ie,
$\semappp{F}{f}{\rho}{v}$ is the meaning of $f$ applied to $v$, when $f$ is a $\lambda$-expression.
Informally, if $f$ is \ensuremath{\lambda_{\Varid{q}}\Varid{x}.(\Varid{e}_{\mathrm{1}})\{\Varid{e}_{\mathrm{2}}\}} then the meaning of this is determined by first evalusing \ensuremath{\Varid{e}_{\mathrm{1}}},
if this succeeds the input is in the domain, and we proceed to evaluate \ensuremath{\Varid{e}_{\mathrm{2}}}.
In more detail, $ws$ is defined by:
\begin{itemize}
\item $\rho' \inb \semapp{X}{\ensuremath{\Varid{e}_{\mathrm{1}}}}{(\rho[x \fmapsto v])}$ \\
  $x$ is the input to the function, we extend the environment with $x$ mapped to the value $v$.
  Then $\semap{X}{e_1}$ will produce all possible environments binding the variables defined in $e_1$.
\item $\ensuremath{\mathbf{let}\;\Varid{ss}} = \semapp{E}{\ensuremath{\Varid{e}_{\mathrm{1}}}}{\rho'}$ \\
  For each of these $\rho'$ we evaluate $e_1$, yielding a sequence.
\item $\squash(\ensuremath{\Varid{ss}}) \neq []$ \\
  If there are no values (\ie, non-empty sets) in the sequence, then $e_1$ failed, and we discard this.
\item $\ensuremath{\mathbf{let}\;\Varid{w}} = \once(\semapp{D}{e_2}{\rho'})$ \\
  If $e_1$ did not fail, get the denotation of $e_2$ in the same environment.
  The {\once} function gets the set out of a singleton sequence, if it is not a singleton then
  $e_2$ failed or is multi-valued (which we do not allow) and so we get the empty set.
\item $[ \ensuremath{\textbf{if}\;\Varid{s}\mathrel{=}}\mempty\,\ensuremath{\textbf{then}}\,\mempty\,\ensuremath{\textbf{else}}\, w\,|\, s \inb ss ]$ \\
  For each non-empty element of the sequence produced by $e_1$ we use the meaning of $e_2$.
\item $\hat{\bigcup}\{\mydots\}$ \\
  The set of sequences from all the different $\rho'$ are combined using $\hat{\bigcup}$;
  this may introduce raggedness.
\end{itemize}
The $ws$ is a sequence of the possible values that the function application can produce.
Since we want the function to behave like a mathematical function, it should only have a single denotation.
The $r$ is the intersection of all the sets in the sequence, this is what all the evaluations of $e_2$ have
in common, so this is what to return.  Finally, per Tim's wish, {\trunc}, will truncate the sequence after
the first non-empty set, replacing it with $r$.

Note that if the domain fails for the given value, $v$, then $semap{F}{}$ will return the empty sequence, $[]$.
Whereas, if the range fails $semap{F}{}$ will return a sequence of empty sets.
The range can fail for several reasons:
\begin{enumerate}
\item the range actually fails
\item there are multiple elements in the sequence whose intersection is empty
\item for a higher order function, the argument failed
\item for a higher order argument, the closedness test failed (added my the \textbf{K} desugaring)
\end{enumerate}

Below I will use \ensuremath{\Varid{a}||\Varid{b}} for unordered choice.  It can easily be defined by \ensuremath{\textbf{block}\;\{\mskip1.5mu \Varid{x}\mathbin{:}\textbf{int};\,\textbf{if}\;(\Varid{x}\mathrel{=}\mathrm{0})\;\{\mskip1.5mu \Varid{a}\mskip1.5mu\}\;\textbf{else}\;\{\mskip1.5mu \Varid{b}\mskip1.5mu\}\mskip1.5mu\}}.

Example 1, \ensuremath{\textbf{fun}_{\Varid{c}}(\mathrm{0})\{\mathrm{5}\}}, desugared \ensuremath{\lambda_{\Varid{c}}\Varid{x}.(\Varid{x}\mathrel{=}\mathrm{0})\{\mathrm{5}\}}, $\semappp{F}{\ensuremath{\lambda_{\Varid{c}}\Varid{x}.(\Varid{x}\mathrel{=}\mathrm{0})\{\mathrm{5}\}}}{\myenv{}}{v}$:

\begin{tabular}{|c|c|c|c|c|c|c|c|c|}
  \hline
  $v$ & $\semapp{X}{e_1}{\myenv{ \mybind{x}{v}}} $ & \mbox{ $ss$ } & $\squash(ss)$ & $w$ & $ \{ \ensuremath{\textbf{if}} \mydots$ \} & $ws$ & $r$ & $\trunc(r,ws)$\\
  \hline
  $0$ & $\{ \myenv{ \mybind{x}{0} } \}$ & \{ $[\{0\}]$ \} & $[\{0\}]$ & \{5\} & \{[\{5\}]\} & [\{5\}] & \{5\} & [\{5\}] \\
  \hline
  $1$ & $\{ \myenv{ \mybind{x}{1} } \}$ & \{ $[\{\}]$ \} & $[]$ & & \{\} & [] & \{\} & [] \\
  \hline
\end{tabular} \\[4mm]

Example 2, \ensuremath{\textbf{fun}_{\Varid{c}}(\mathrm{0}\hspace{0.19em}\mathop{\rule[-0.1em]{0.15em}{0.75em}}\hspace{0.2em}\mathrm{1})\{\mathrm{5}\}}, desugared \ensuremath{\lambda_{\Varid{c}}\Varid{x}.((\Varid{x}\mathrel{=}\mathrm{0})\hspace{0.19em}\mathop{\rule[-0.1em]{0.15em}{0.75em}}\hspace{0.2em}(\Varid{x}\mathrel{=}\mathrm{1}))\{\mathrm{5}\}}, $\semappp{F}{\ensuremath{\lambda_{\Varid{c}}\Varid{x}.((\Varid{x}\mathrel{=}\mathrm{0})\hspace{0.19em}\mathop{\rule[-0.1em]{0.15em}{0.75em}}\hspace{0.2em}(\Varid{x}\mathrel{=}\mathrm{1}))\{\mathrm{5}\}}}{\myenv{}}{v}$:

\begin{tabular}{|c|c|c|c|c|c|c|c|c|}
  \hline
  $v$ & $\semapp{X}{e_1}{\myenv{ \mybind{x}{v}}} $ & \mbox{ $ss$ } & $\squash(ss)$ & $w$ & $ \{ \ensuremath{\textbf{if}} \mydots$ \} & $ws$ & $r$ & $\trunc(r,ws)$\\
  \hline
  $0$ & $\{ \myenv{ \mybind{x}{0} } \}$ & \{ $[\{0\},\{\}]$ \} & $[\{0\}]$ & \{5\} & \{[\{5\},\{\}]\} & [\{5\},\{\}] & \{5\} & [\{5\}] \\
  \hline
  $1$ & $\{ \myenv{ \mybind{x}{1} } \}$ & \{ $[\{\}.\{1\}]$ \} & $[\{1\}]$ & \{5\} & \{[\{\},\{5\}]\} & [\{\},\{5\}] & \{5\} & [\{\},\{5\}] \\
  \hline
  $2$ & $\{ \myenv{ \mybind{x}{2} } \}$ & \{ $[\{\},\{\}]$ \} & $[]$ & & \{\} & [] & \{\} & [] \\
  \hline
\end{tabular} \\[4mm]

Example 3, \ensuremath{\textbf{fun}_{\Varid{c}}(\mathrm{0}\hspace{0.19em}\mathop{\rule[-0.1em]{0.15em}{0.75em}}\hspace{0.2em}\mathrm{1}\hspace{0.19em}\mathop{\rule[-0.1em]{0.15em}{0.75em}}\hspace{0.2em}\mathrm{0})\{\mathrm{5}\}}, desugared \ensuremath{\lambda_{\Varid{c}}\Varid{x}.((\Varid{x}\mathrel{=}\mathrm{0})\hspace{0.19em}\mathop{\rule[-0.1em]{0.15em}{0.75em}}\hspace{0.2em}(\Varid{x}\mathrel{=}\mathrm{1})\hspace{0.19em}\mathop{\rule[-0.1em]{0.15em}{0.75em}}\hspace{0.2em}\mathrm{0})\{\mathrm{5}\}}, $\semappp{F}{\ensuremath{\lambda_{\Varid{c}}\Varid{x}.((\Varid{x}\mathrel{=}\mathrm{0})\hspace{0.19em}\mathop{\rule[-0.1em]{0.15em}{0.75em}}\hspace{0.2em}(\Varid{x}\mathrel{=}\mathrm{1})\hspace{0.19em}\mathop{\rule[-0.1em]{0.15em}{0.75em}}\hspace{0.2em}\mathrm{0})\{\mathrm{5}\}}}{\myenv{}}{v}$:

\begin{tabular}{|c|c|c|c|c|c|c|c|c|}
  \hline
  $v$ & $\semapp{X}{e_1}{\myenv{ \mybind{x}{v}}} $ & \mbox{ $ss$ } & $\squash(ss)$ & $w$ & $ \{ \ensuremath{\textbf{if}} \mydots$ \} & $ws$ & $r$ & $\trunc(r,ws)$\\
  \hline
  $0$ & $\{ \myenv{ \mybind{x}{0} } \}$ & \{ $[\{0\},\{\},\{0\}]$ \} & $[\{0\},\{0\}]$ & \{5\} & \{[\{5\},\{\},\{5\}]\} & [\{5\},\{\},\{5\}] & \{5\} & [\{5\}] \\
  \hline
  $1$ & $\{ \myenv{ \mybind{x}{1} } \}$ & \{ $[\{\}.\{1\},\{\}]$ \} & $[\{1\}]$ & \{5\} & \{[\{\},\{5\},\{\}]\} & [\{\},\{5\},\{\}] & \{5\} & [\{\},\{5\}] \\
  \hline
  $2$ & $\{ \myenv{ \mybind{x}{2} } \}$ & \{ $[\{\},\{\},\{\}]$ \} & $[]$ & & \{\} & [] & \{\} & [] \\
  \hline
\end{tabular} \\[4mm]

Example 4, \ensuremath{\textbf{fun}_{\Varid{c}}(\mathrm{0}||\mathrm{1})\{\mathrm{5}||\mathrm{6}\}}, desugared \ensuremath{\lambda_{\Varid{c}}\Varid{x}.((\Varid{x}\mathrel{=}\mathrm{0})||(\Varid{x}\mathrel{=}\mathrm{1})\;(\mathrm{5}||\mathrm{6}))\{\cdot \}}, $\semappp{F}{\ensuremath{\lambda_{\Varid{c}}\Varid{x}.((\Varid{x}\mathrel{=}\mathrm{0})||(\Varid{x}\mathrel{=}\mathrm{1}))\{\mathrm{5}||\mathrm{6}\}}}{\myenv{}}{v}$:

\begin{tabular}{|c|c|c|c|c|c|c|c|c|}
  \hline
  $v$ & $\semapp{X}{e_1}{\myenv{ \mybind{x}{v}}} $ & \mbox{ $ss$ } & $\squash(ss)$ & $w$ & $ \{ \ensuremath{\textbf{if}} \mydots$ \} & $ws$ & $r$ & $\trunc(r,ws)$\\
  \hline
  $0$ & $\{ \myenv{ \mybind{x}{0} } \}$ & \{ $[\{0\}]$ \} & $[\{0\}]$ & \{5,6\} & \{[\{5,6\}]\} & [\{5,6\}] & \{5,6\} & [\{5,6\}] \\
  \hline
  $1$ & $\{ \myenv{ \mybind{x}{1} } \}$ & \{ $[\{1\}]$ \} & $[\{1\}]$ & \{5,6\} & \{[\{5,6\}]\} & [\{5,6\}] & \{5,6\} & [\{5,6\}] \\
  \hline
  $2$ & $\{ \myenv{ \mybind{x}{2} } \}$ & \{ $[\{\}]$ \} & $[]$ & & \{\} & [] & \{\} & [] \\
  \hline
\end{tabular} \\[4mm]

To obtain the denotation of a function the $\semap{E}{}$ function will select all functions
that are compatible with the function definition.
\begin{itemize}
\item $\{ f \,|\, f \inb [\semfcn]$ \\
  For all functions (\ie, list of mathematical functions),
\item $\forall v \inb \Value .$ \\
  for all values $v$ (\ie, possible function arguments)
\item $\seqcompat(\applyo(f,v), \semappp{F}{\ensuremath{\lambda_{\Varid{q}}\Varid{x}.(\Varid{e}_{\mathrm{1}})\{\Varid{e}_{\mathrm{2}}\}}}{\rho}{v})$ \\
  the result of applying the function, $f$, to the argument $v$ is compatible
  with the applying the function definition to the argument.
\item $\semopen{q}{(fs)}$ \\
  Finally, if the is an open function get the set of all functions with a larger domain.
\end{itemize}

The result of $\applyo(f,v)$ will be a sequence of sets.
Each set is either empty or a singleton set, because the sets come from applying
mathematical functions to a single value.
On the other hand, the result of $\semappp{F}{\ensuremath{\lambda_{\Varid{q}}\Varid{x}.(\Varid{e}_{\mathrm{1}})\{\Varid{e}_{\mathrm{2}}\}}}{\rho}{v}$ is a sequence
of sets where the sets can contain multiple values.  This is because the function definition may
denote many functions.
These two are compatible (checked by \seqcompat) if the sequences have the same length and the
individual sets are compatible (checked by \setcompat).
A set from $\applyo(f,v)$ is compatible with a set from $\semap{F}{\mydots}$ if they are
both empty or, for non-empty sets, if the first is a subset of the second.

\lennart{I need to think more about decides functions.  This doesn't smell quite right.}

%  To obtain the denotation of a function the $\semap{E}{}$ function will iterate through all possible
%  value and use $\semap{F}{}$ to find out how the function behaves for each of them.  It will then use this
%  to stitch together (the singleton sequence of) the set of functions that is the denotation.  In more detail:
%  \begin{itemize}
%  \item $v \inb \Value$ \\
%    For all values $v$,
%  \item $|let r| = \semapp{F}{|lamOC q x e_1 e_2|}{\rho}\;v$ \\
%    find the function result (a sequence of sets) for that $v$,
%  \item $r \neq []$ \\
%    and only keep the ones where the function succeeds.
%  \item $[ \dist(v, s) \,||\, s \inb r ]$ \\
%    Pair up the argument with each possible result.
%    The pairs are formed by $\mapsto$ to indicate that they will form a function.
%  \item $\combine(|fs|)$ \\
%    Form the actual closed functions from all the pieces.
%    \combine is described in more detail below.
%  \item $\semopen{q}{(\mydots)}$ \\
%    If the function is open, select all functions that are compatible
%    with result of \combine.
%  \item \succeeds \\
%    Check that all sequences in |fs| contain a non-empty set.  If not, it means that the range failed when the domain succeeded.
%    Omit this check for a \textbf{decides} function.
%  \end{itemize}
%  
%  Continuing with the examples.
%  
%  Example 1, |functionOC c 0 5|
%  
%  \begin{tabular}{||c||c||c||c||c||}
%    \hline
%    $v$ & $r$ & $r \neq []$ & $[ \dist(v, s) \,||\, s \inb r ]$ & |fs| \\
%    \hline
%    $0$ & $[\{5\}]$ & $[\{5\}]$ & $[\{0 \mapsto 5\}]$ & \multirow{2}{*}{$\{[\{0 \mapsto 5\}]\}$}\\
%    \cline{1-4}
%    $1$ & $[]$ & & & \\
%    \hline
%  \end{tabular}\\
%  \begin{tabular}{||c||c||c||c||}
%    \hline
%    |fs| & \group & \maplist(\cross) & \sequence \\
%    \hline
%    $\{[\{0 \mapsto 5\}]\}$ & $[\{0 \mapsto 5\}]$ & $[\{0 \mapsto 5\}]$ & $\{[0 \mapsto 5]\}$ \\
%    \hline
%  \end{tabular}\\[5mm]
%  
%  Example 2, |functionOC c (0 `choice` 1) 5|
%  \lennart{WRONG}
%  \begin{tabular}{||c||c||c||c||c||}
%    \hline
%    $v$ & $r$ & $r \neq []$ & $[ \dist(v, s) \,||\, s \inb r ]$ & |fs| \\
%    \hline
%    $0$ & $[\{5\}]$ & $[\{5\}]$ & $[\{0 \mapsto 5\}]$ & \multirow{3}{*}{$\{[\{0 \mapsto 5\}],[\{\},\{1 \mapsto 5\}]\}$}\\
%    \cline{1-4}
%    $1$ & $[\{\},\{5\}]$ & $[\{\},\{5\}]$ & $[\{\},\{1 \mapsto 5\}]$ & \\
%    \cline{1-4}
%    $2$ & $[]$ & & & \\
%    \hline
%  \end{tabular}\\
%  \begin{tabular}{||c||c||c||c||}
%    \hline
%    |fs| & \group & \maplist(\cross) & \sequence \\
%    \hline
%    $\{[\{0 \mapsto 5\}],[\{\},\{1 \mapsto 5\}]\}$ &
%      $[\{0 \mapsto 5\}, \{1 \mapsto 5\}]$ &
%      $[\{0 \mapsto 5\}, \{1 \mapsto 5\}]$ &
%      TBD \\
%    \hline
%  \end{tabular}\\[5mm]
%  
%  Example 3, same as example 2
%  
%  Example 4, |functionOC c (0 `uchoice` 1) (5 `uchoice` 6)|
%  \lennart{WRONG}
%  
%  \begin{tabular}{||c||c||c||c||c||}
%    \hline
%    $v$ & $r$ & $r \neq []$ & $[ \dist(v, s) \,||\, s \inb r ]$ & |fs| \\
%    \hline
%    $0$ & $[\{5,6\}]$ & $[\{5,6\}]$ & $[\{0 \mapsto 5, 0 \mapsto 6\}]$ & \multirow{3}{*}{
%      \makecell{$\{[\{0 \mapsto 5, 0 \mapsto 6,$ \\
%                $1 \mapsto 5, 1 \mapsto 6\}]\}$}
%    }\\
%    \cline{1-4}
%    $1$ & $[\{5,6\}]$ & $[\{5,6\}]$ & $[\{1 \mapsto 5, 0 \mapsto 6\}]$ & \\
%    \cline{1-4}
%    $2$ & $[]$ & & & \\
%    \hline
%  \end{tabular}\\
%  \begin{tabular}{||c||c||c||c||}
%    \hline
%    |fs| & \group & \maplist(\cross) & \sequence \\
%    \hline
%      \makecell{$\{[\{0 \mapsto 5, 0 \mapsto 6,$ \\
%                $1 \mapsto 5, 1 \mapsto 6\}]\}$} &
%      \makecell{$[\{0 \mapsto 5, 0 \mapsto 6,$ \\
%          $1 \mapsto 5, 1 \mapsto 6\}]$} &
%      TBD & TBD \\
%    \hline
%  \end{tabular}\\[5mm]
%  
%  
